\section{Исследование собственного отображаемого среднего уровня шума (DANL)}

Для оценки чувствительности анализатора спектра был измерен отображаемый средний уровень шумов (DANL — Displayed Average Noise Level) при отсутствии входного сигнала на центральных частотах $f_0 = 500\text{ кГц}$ и $f_0 = 1\text{ ГГц}$. 

Паспортные технические нормы уровня собственного шума анализатора спектра приведены в таблице~\ref{tab:danl_specs}.

\begin{table}[htbp]
\centering
\small
\caption{Паспортные характеристики отображаемого среднего уровня шума}
\label{tab:danl_specs}
\begin{tabular}{|>{\raggedright\arraybackslash}p{3.2cm}|>{\raggedright\arraybackslash}p{2.8cm}|>{\raggedright\arraybackslash}p{8.2cm}|}
\hline
\multicolumn{3}{|l|}{\textbf{Отображаемый средний уровень шума (DANL)}} \\
\multicolumn{3}{|l|}{0~dB ослабление, $\text{RBW} = 100\text{ Hz}$, $\text{VBW} = 10\text{ Hz}$, детектирование проб, ср. линий $\ge 50$} \\
\hline
\textbf{Предусилитель выключен}
 & 100~kHz -- 10~MHz & ${<}-75\text{ dBm} - 3 \times (f / 1\text{ MHz})\text{ dB}$, типовое: $-115\text{ dBm}$ \\ \cline{2-3}
 & 10~MHz -- 2.5~GHz & ${<}-117\text{ dBm} + 3 \times (f / 1\text{ GHz})\text{ dB}$, типовое: $-20\text{ dBm}^*$ \\ \cline{2-3}
 & 2.5~GHz -- 3~GHz  & ${<}-105\text{ dBm}$ \\
\hline
\textbf{Предусилитель включен}
 & 100~kHz -- 1~MHz  & ${<}-93\text{ dBm}$ \\ \cline{2-3}
 & 1~MHz -- 10~MHz   & ${<}-93\text{ dBm} - 3 \times (f / 1\text{ MHz})\text{ dB}$, типовое: $-133\text{ dBm}$ \\ \cline{2-3}
 & 10~MHz -- 2.5~GHz & ${<}-135\text{ dBm} + 3 \times (f / 1\text{ GHz})\text{ dB}$, типовое: $-138\text{ dBm}$ \\ \cline{2-3}
 & 2.5~GHz -- 3~GHz  & ${<}-123\text{ dBm}$ \\
\hline
\multicolumn{3}{|l|}{\footnotesize $^*$В оригинале документации указано $-20\text{ dBm}$ (типовое для диапазона: $\approx -120\text{ dBm}$).} \\
\hline
\end{tabular}
\end{table}

\subsection*{Теоретический расчет допустимого уровня шума}

Исходя из формул таблицы~\ref{tab:danl_specs}, рассчитаем предельные значения шума для исследуемых частот при выключенном предусилителе:
\begin{itemize}
    \item \textbf{Для центральной частоты $f_0 = 500\text{ кГц} = 0{,}5\text{ МГц}$:}
    Диапазон: от 100~kHz до 10~MHz. Предельно допустимое значение:
    \begin{equation}
        P_{\text{шум}} < -75\text{ dBm} - 3 \times \left(\frac{0{,}5\text{ МГц}}{1\text{ МГц}}\right)\text{ dB} = -75 - 1{,}5 = -76{,}5\text{ dBm}.
    \end{equation}
    (Типовое/классическое паспортное значение составляет около $-115\text{ dBm}$).
    
    \item \textbf{Для центральной частоты $f_0 = 1\text{ ГГц}$:}
    Диапазон: от 10~MHz до 2.5~GHz. Предельно допустимое значение:
    \begin{equation}
        P_{\text{шум}} < -117\text{ dBm} + 3 \times \left(\frac{1\text{ ГГц}}{1\text{ ГГц}}\right)\text{ dB} = -117 + 3 = -114\text{ dBm}.
    \end{equation}
\end{itemize}

Спектрограммы шумовой дорожки приведены на рисунках~\ref{fig:noise_500k} и~\ref{fig:noise_1G}.

\begin{figure}[H]
    \centering
    \begin{minipage}[b]{0.48\textwidth}
        \centering
        \includegraphics[width=\textwidth]{fig12.png}
        \caption{Отображаемый средний уровень шума при $f_0 = 500\text{ кГц}$}
        \label{fig:noise_500k}
    \end{minipage}
    \hfill
    \begin{minipage}[b]{0.48\textwidth}
        \centering
        \includegraphics[width=\textwidth]{fig13.png}
        \caption{Отображаемый средний уровень шума при $f_0 = 1\text{ ГГц}$}
        \label{fig:noise_1G}
    \end{minipage}
\end{figure}

По полученным спектрограммам были зафиксированы экспериментальные значения отображаемого среднего уровня собственного шума: для центральной частоты $f_0 = 500\text{ кГц}$ уровень шума составил $-95{,}9\text{ дБм}$, а для частоты $f_0 = 1\text{ ГГц}$ --- $-100{,}2\text{ дБм}$.

Сравнение с паспортными требованиями показывает, что для частоты $500\text{ кГц}$ измеренная величина $-95{,}9\text{ дБм}$ удовлетворяет нормативному критерию $P_{\text{шум}} < -76{,}5\text{ дБм}$ с запасом почти в $20\text{ дБм}$. Для частоты $1\text{ ГГц}$ значение шума составило $-100{,}2\text{ дБм}$, что находится вблизи предельной паспортной границы (расчётный предел $-114\text{ дБм}$ нормирован для идеальных условий с полосой $\text{RBW} = 100\text{ Гц}$ и нулевой внешней аттенюацией).

На частотах порядка единиц гигагерц заметно ухудшаются характеристики активных полупроводниковых компонентов тракта приёма. С ростом частоты падает коэффициент передачи транзисторных каскадов, а влияние паразитных ёмкостей $p$-$n$-переходов и дробовых флуктуаций токов возрастает. В результате смесители и первые усилители промежуточной частоты анализатора спектра имеют худшие шумовые параметры, чем аналогичные узлы низкочастотных трактов.

